% =============================================================================
% MLSys·im Tutorial — Module 2: Advanced Single-Node Analysis
% =============================================================================
\documentclass[aspectratio=169, 12pt]{beamer}
\usepackage{../../../slides/assets/beamerthememlsys}

\mlsyssetup{
  volume       = {Tutorial},
  chapter      = {Module 2},
  logo         = {../../../slides/assets/img/logo-mlsysbook.png},
  instlogo     = {../../../slides/assets/img/logo-harvard.png},
  chaptertitle = {MLSys·im: Advanced Single-Node Analysis},
}

% --- Fonts & Packages ---
\usepackage[T1]{fontenc}
\usepackage[scaled=0.9]{helvet}
\usepackage{courier}
\renewcommand{\familydefault}{\sfdefault}
\usepackage{amsmath}
\usepackage{booktabs}
\usepackage[table]{xcolor}
\usepackage{listings}
\usepackage{tikz}

% --- Code listings ---
\lstset{
  language=Python,
  basicstyle=\ttfamily\footnotesize,
  keywordstyle=\color{crimson}\bfseries,
  stringstyle=\color{datastroke},
  commentstyle=\color{midgray}\itshape,
  backgroundcolor=\color{computeblue!20},
  frame=single,
  rulecolor=\color{computestroke},
  numbers=none,
  breaklines=true,
  columns=fullflexible,
  keepspaces=true,
  showstringspaces=false,
  xleftmargin=4pt,
  xrightmargin=4pt,
  aboveskip=6pt,
  belowskip=4pt,
}

\newcommand{\mlsysim}{\texttt{mlsysim}}
\graphicspath{{images/}}

\title{MLSys·im Tutorial --- Module 2}
\subtitle{Advanced Single-Node Analysis}
\author{Vijay Janapa Reddi}
\institute{Harvard University}
\date{Conference Tutorial}

\begin{document}

\begin{frame}[plain]
  \titlepage
\end{frame}

\begin{frame}{Roadmap: Conference Tutorial}
\centering\small
\begin{tabular}{rll}
\toprule
\textbf{Module} & \textbf{Topic} & \textbf{Status} \\
\midrule
Module 1 & Foundations \& Architecture & \checkmark Done \\
\rowcolor{crimson!12}
\textbf{Module 2} & \textbf{Advanced Single-Node Analysis} & \textbf{$\leftarrow$ You are here} \\
Module 3 & Scale, Dollars, and Carbon & \\
Module 4 & Design Space Exploration \& Synthesis & \\
\bottomrule
\end{tabular}
\end{frame}

\section{The Data Wall}

\begin{frame}{Beyond FLOPs: The Data Wall}
GPUs are so fast they often starve waiting for data.
\begin{itemize}
  \item \textbf{Ingestion (I/O):} Can the storage sub-system (NVMe, PCIe) push bytes fast enough?
  \item \textbf{Transformation (CPU):} Can the CPUs decode JPEGs and run augmentations fast enough to keep the GPUs busy?
\end{itemize}
\end{frame}

\begin{frame}[fragile]{Live Demo: Uncovering the Data Wall}
\note{This corresponds to Scenario C in the Code Cookbook}
\begin{lstlisting}[language=Python]
from mlsysim.solvers import DataModel, TransformationModel
from mlsysim.hardware.registry import Hardware
from mlsysim.core.constants import Q_

demand_rate = Q_("40000 1/s") * Q_("150 KB") # ~6 GB/s

# 1. Check if the DGX A100 PCIe bus can handle the bandwidth
data_result = DataModel().solve(
    workload_data_rate=demand_rate, hardware=Hardware.Cloud.A100)
print(f"I/O Stalled: {data_result.is_stalled}") # False

# 2. Check if the CPUs can decode/augment images fast enough
transform_result = TransformationModel().solve(
    batch_size=40000,
    cpu_throughput_per_worker_hz=850, # JPEG decode + crop
    num_workers=64                    # 8 CPUs per GPU
)
print(f"CPU Stalled: {transform_result.is_bottleneck}") # True
\end{lstlisting}
\end{frame}

\section{LLM Serving \& Memory Walls}

\begin{frame}{The Two Phases of LLM Serving}
LLM inference is not a single mathematical operation; it is a stateful process with two distinct physical regimes.
\vfill
\begin{columns}[T]
\column{0.5\textwidth}
\textbf{1. Pre-fill (Prompt Processing)}
\begin{itemize}
  \item Processes all prompt tokens in parallel.
  \item Matrix-Matrix multiplication.
  \item High arithmetic intensity $\Rightarrow$ \textbf{Compute Bound}.
  \item Metric: Time-to-First-Token (TTFT).
\end{itemize}

\column{0.5\textwidth}
\textbf{2. Decoding (Token Generation)}
\begin{itemize}
  \item Generates one token at a time autoregressively.
  \item Matrix-Vector multiplication.
  \item Low arithmetic intensity $\Rightarrow$ \textbf{Memory Bound}.
  \item Metric: Inter-Token Latency (ITL).
\end{itemize}
\end{columns}
\end{frame}

\section{Algorithmic Optimizations}

\begin{frame}[fragile]{Speculative Decoding (Scenario E)}
\emph{What if we use a smaller, cheaper model to "guess" tokens, and the large model just verifies them?}
\begin{lstlisting}[language=Python]
from mlsysim.solvers import ServingModel
from mlsysim.hardware.registry import Hardware
from mlsysim.models.registry import Models

target_model = Models.Language.Llama3_70B
draft_model = Models.Language.Llama3_8B
hardware = Hardware.Cloud.H100

solver = ServingModel()
base_result = solver.solve(target_model, hardware, seq_len=2048, batch_size=1)
print(f"Standard ITL: {base_result.itl.to('ms'):.2f}")

spec_result = solver.solve(
    target_model, hardware, seq_len=2048, batch_size=1,
    draft_model=draft_model, draft_acceptance_rate=0.75)
print(f"Speedup: {base_result.itl / spec_result.itl:.2f}x")
\end{lstlisting}
\end{frame}

\section{The Reasoning Wall}

\begin{frame}[fragile]{Wall 12: Inference-Time Compute (The Reasoning Wall)}
With models like OpenAI o1, compute scaling shifts from training to inference. The model generates $K$ hidden reasoning steps before answering.
\begin{lstlisting}[language=Python]
from mlsysim.engine.solver import InferenceScalingModel

reasoning_solver = InferenceScalingModel()
result = reasoning_solver.solve(
    model=Models.Language.Llama3_8B, hardware=Hardware.Cloud.H100,
    reasoning_steps=128,  # K hidden CoT steps
    precision="fp16", efficiency=0.5)

print(f"Reasoning Time: {result.total_reasoning_time.to('s'):.1f}")
print(f"Energy per Query: {result.energy_per_query.to('J'):.1f}")
\end{lstlisting}
\vfill
\textbf{Impact:} Inference shifts back toward being \emph{compute-bound} as generation length dominates prefill.
\end{frame}

\begin{frame}{Summary: Module 2}
\begin{enumerate}
  \item \textbf{Data Wall:} Real-world throughput is often bottlenecked by CPU transformation, not GPU FLOPs.
  \item \textbf{Serving:} Pre-fill is Compute-bound; Decoding is Memory-bound.
  \item \textbf{Algorithmic Optimizations:} We can instantly model the speedup of Speculative Decoding and Quantization.
  \item \textbf{Reasoning Wall:} Hidden CoT tokens push inference back into the compute-bound regime.
\end{enumerate}
\vspace{1em}
\begin{center}
\textit{Next up: Module 3 - Scale, Dollars, and Carbon!}
\end{center}
\end{frame}

\end{document}